\section{Quasi-coherent modules}
\label{section-quasi-coherent}

\noindent
In this section we introduce an abstract notion of
quasi-coherent $\mathcal{O}_X$-module. This notion is very
useful in algebraic geometry, since quasi-coherent modules
on a scheme have a good description on any affine open.
However, we warn the reader that
in the general setting of (locally) ringed spaces
this notion is not well behaved at all. The category of
quasi-coherent sheaves is not abelian in general, infinite
direct sums of quasi-coherent sheaves aren't quasi-coherent, etc, etc.

\begin{definition}
\label{definition-quasi-coherent}
Let $(X, \mathcal{O}_X)$ be a ringed space.
Let $\mathcal{F}$ be a sheaf of $\mathcal{O}_X$-modules.
We say that $\mathcal{F}$ is a {\it quasi-coherent
sheaf of $\mathcal{O}_X$-modules} if for every
point $x \in X$ there exists an open neighbourhood
$x\in U \subset X$ such that $\mathcal{F}|_U$
is isomorphic to the cokernel of a map
$$
\bigoplus\nolimits_{j \in J}
\mathcal{O}_U
\longrightarrow
\bigoplus\nolimits_{i \in I}
\mathcal{O}_U
$$
The category of quasi-coherent $\mathcal{O}_X$-modules
is denoted $\QCoh(\mathcal{O}_X)$.
\end{definition}

\noindent
The definition means that $X$ is covered by open sets $U$
such that $\mathcal{F}|_U$ has a {\it presentation}
of the form
$$
\bigoplus\nolimits_{j \in J}
\mathcal{O}_U
\longrightarrow
\bigoplus\nolimits_{i \in I}
\mathcal{O}_U
\longrightarrow
\mathcal{F}|_U
\longrightarrow
0.
$$
Here presentation signifies that the displayed
sequence is exact. In other words
\begin{enumerate}
\item for every point $x$ of $X$ there exists
an open neighbourhood such that $\mathcal{F}|_U$
is generated by global sections, and
\item for a suitable choice of these sections
the kernel of the associated surjection is also
generated by global sections.
\end{enumerate}

\begin{lemma}
\label{lemma-direct-sum-quasi-coherent}
Let $(X, \mathcal{O}_X)$ be a ringed space.
The direct sum of two quasi-coherent $\mathcal{O}_X$-modules is
a quasi-coherent $\mathcal{O}_X$-module.
\end{lemma}

\begin{proof}
Omitted.
\end{proof}

\begin{remark}
\label{remark-infinite-direct-sum-quasi-coherent-not}
Warning: It is not true in general that an infinite
direct sum of quasi-coherent $\mathcal{O}_X$-modules
is quasi-coherent. For more esoteric behaviour of quasi-coherent
modules see Example \ref{example-quasi-coherent}.
\end{remark}

\begin{lemma}
\label{lemma-pullback-quasi-coherent}
Let $f : (X, \mathcal{O}_X) \to (Y, \mathcal{O}_Y)$
be a morphism of ringed spaces.
The pullback $f^*\mathcal{G}$ of a quasi-coherent
$\mathcal{O}_Y$-module is quasi-coherent.
\end{lemma}

\begin{proof}
Arguing as in the proof of Lemma \ref{lemma-pullback-locally-generated}
we may assume $\mathcal{G}$ has a global presentation by
direct sums of copies of $\mathcal{O}_Y$.
We have seen that $f^*$ commutes with all colimits,
and is right exact, see Lemma \ref{lemma-exactness-pushforward-pullback}.
Thus if we have an exact sequence
$$
\bigoplus\nolimits_{j \in J}
\mathcal{O}_Y
\longrightarrow
\bigoplus\nolimits_{i \in I}
\mathcal{O}_Y
\longrightarrow
\mathcal{G}
\longrightarrow
0
$$
then upon applying $f^*$ we obtain the exact sequence
$$
\bigoplus\nolimits_{j \in J}
\mathcal{O}_X
\longrightarrow
\bigoplus\nolimits_{i \in I}
\mathcal{O}_X
\longrightarrow
f^*\mathcal{G}
\longrightarrow
0.
$$
This implies the lemma.
\end{proof}

\noindent
This gives plenty of examples of quasi-coherent sheaves.

\begin{lemma}
\label{lemma-construct-quasi-coherent-sheaves}
Let $(X, \mathcal{O}_X)$ be ringed space.
Let $\alpha : R \to \Gamma(X, \mathcal{O}_X)$ be a ring homomorphism from
a ring $R$ into the ring of global sections on $X$.
Let $M$ be an $R$-module.
The following three constructions give canonically isomorphic
sheaves of $\mathcal{O}_X$-modules:
\begin{enumerate}
\item Let $\pi : (X, \mathcal{O}_X) \longrightarrow (\{*\}, R)$
be the morphism of ringed spaces with $\pi : X \to \{*\}$
the unique map and with $\pi$-map $\pi^\sharp$ the given map
$\alpha : R \to \Gamma(X, \mathcal{O}_X)$. Set $\mathcal{F}_1 = \pi^*M$.
\item Choose a presentation
$\bigoplus_{j \in J} R \to \bigoplus_{i \in I} R \to M \to 0$.
Set
$$
\mathcal{F}_2 = \Coker\left(
\bigoplus\nolimits_{j \in J} \mathcal{O}_X
\to
\bigoplus\nolimits_{i \in I} \mathcal{O}_X
\right).
$$
Here the map on the component $\mathcal{O}_X$ corresponding to $j \in J$
given by the section $\sum_i \alpha(r_{ij})$ where the $r_{ij}$
are the matrix coefficients of the map in the presentation of $M$.
\item Set $\mathcal{F}_3$ equal to the sheaf associated to the presheaf
$U \mapsto \mathcal{O}_X(U) \otimes_R M$, where the map
$R \to \mathcal{O}_X(U)$ is the composition of $\alpha$ and
the restriction map $\mathcal{O}_X(X) \to \mathcal{O}_X(U)$.
\end{enumerate}
This construction has the following properties:
\begin{enumerate}
\item The resulting sheaf of $\mathcal{O}_X$-modules
$\mathcal{F}_M = \mathcal{F}_1 = \mathcal{F}_2 = \mathcal{F}_3$
is quasi-coherent.
\item The construction gives a functor from
the category of $R$-modules to the category of quasi-coherent
sheaves on $X$ which commutes with arbitrary colimits.
\item For any $x \in X$ we have
$\mathcal{F}_{M, x} = \mathcal{O}_{X, x} \otimes_R M$
functorial in $M$.
\item Given any $\mathcal{O}_X$-module
$\mathcal{G}$ we have
$$
\Mor_{\mathcal{O}_X}(\mathcal{F}_M, \mathcal{G})
=
\Hom_R(M, \Gamma(X, \mathcal{G}))
$$
where the $R$-module structure on $\Gamma(X, \mathcal{G})$
comes from the $\Gamma(X, \mathcal{O}_X)$-module structure via
$\alpha$.
\end{enumerate}
\end{lemma}

\begin{proof}
The isomorphism between $\mathcal{F}_1$ and $\mathcal{F}_3$
comes from the fact that $\pi^*$ is defined as the sheafification
of the presheaf in (3), see Sheaves, Section
\ref{sheaves-section-ringed-spaces-functoriality-modules}.
The isomorphism between the constructions in (2) and (1) comes
from the fact that the functor $\pi^*$ is right exact, so
$\pi^*(\bigoplus_{j \in J} R) \to \pi^*(\bigoplus_{i \in I} R) \to
\pi^*M \to 0$ is exact, $\pi^*$ commutes with arbitrary
direct sums, see Lemma \ref{lemma-exactness-pushforward-pullback},
and finally the fact that $\pi^*(R) = \mathcal{O}_X$.

\medskip\noindent
Assertion (1) is clear from construction (2).
Assertion (2) is clear since $\pi^*$ has these properties.
Assertion (3) follows from the description of stalks of
pullback sheaves, see
Sheaves, Lemma \ref{sheaves-lemma-stalk-pullback-modules}.
Assertion (4) follows from adjointness of $\pi_*$ and
$\pi^*$.
\end{proof}

\begin{definition}
\label{definition-sheaf-associated}
In the situation of Lemma \ref{lemma-construct-quasi-coherent-sheaves}
we say $\mathcal{F}_M$ is the {\it sheaf associated to the module $M$
and the ring map $\alpha$}. If $R = \Gamma(X, \mathcal{O}_X)$
and $\alpha = \text{id}_R$ we simply say $\mathcal{F}_M$ is the
{\it sheaf associated to the module $M$}.
\end{definition}


\begin{lemma}
\label{lemma-restrict-quasi-coherent}
Let $(X, \mathcal{O}_X)$ be a ringed space.
Set $R = \Gamma(X, \mathcal{O}_X)$.
Let $M$ be an $R$-module.
Let $\mathcal{F}_M$ be the quasi-coherent sheaf of
$\mathcal{O}_X$-modules associated to $M$.
If $g : (Y, \mathcal{O}_Y) \to (X, \mathcal{O}_X)$
is a morphism of ringed spaces, then
$g^*\mathcal{F}_M$ is the sheaf associated
to the $\Gamma(Y, \mathcal{O}_Y)$-module
$\Gamma(Y, \mathcal{O}_Y) \otimes_R M$.
\end{lemma}

\begin{proof}
The assertion follows from the first description
of $\mathcal{F}_M$ in Lemma \ref{lemma-construct-quasi-coherent-sheaves}
as $\pi^*M$, and the following commutative diagram
of ringed spaces
$$
\xymatrix{
(Y, \mathcal{O}_Y) \ar[r]_-\pi \ar[d]_g &
(\{*\}, \Gamma(Y, \mathcal{O}_Y)) \ar[d]^{\text{induced by }g^\sharp} \\
(X, \mathcal{O}_X) \ar[r]^-\pi &
(\{*\}, \Gamma(X, \mathcal{O}_X))
}
$$
(Also use Sheaves, Lemma \ref{sheaves-lemma-push-pull-composition-modules}.)
\end{proof}

\begin{lemma}
\label{lemma-quasi-coherent-module}
Let $(X, \mathcal{O}_X)$ be a ringed space.
Let $x \in X$ be a point.
Assume that $x$ has a fundamental system of quasi-compact neighbourhoods.
Consider any quasi-coherent $\mathcal{O}_X$-module $\mathcal{F}$.
Then there exists an open neighbourhood $U$ of $x$
such that $\mathcal{F}|_U$ is isomorphic to the
sheaf of modules $\mathcal{F}_M$ on $(U, \mathcal{O}_U)$
associated to some $\Gamma(U, \mathcal{O}_U)$-module $M$.
\end{lemma}

\begin{proof}
First we may replace $X$ by an open neighbourhood of $x$
and assume that $\mathcal{F}$ is isomorphic to the
cokernel of a map
$$
\Psi :
\bigoplus\nolimits_{j \in J}
\mathcal{O}_X
\longrightarrow
\bigoplus\nolimits_{i \in I}
\mathcal{O}_X.
$$
The problem is that this map may not be given by
a ``matrix'', because the module of global sections of a direct sum
is in general different from the direct sum of the modules
of global sections.

\medskip\noindent
Let $x \in E \subset X$ be a quasi-compact
neighbourhood of $x$ (note: $E$ may not be open).
Let $x \in U \subset E$ be an open neighbourhood of
$x$ contained in $E$.
Next, we proceed as in the proof of
Lemma \ref{lemma-section-direct-sum-quasi-compact}.
For each $j \in J$ denote
$s_j \in \Gamma(X, \bigoplus\nolimits_{i \in I} \mathcal{O}_X)$
the image of the section $1$ in the summand $\mathcal{O}_X$
corresponding to $j$. There exists a finite collection of opens
$U_{jk}$, $k \in K_j$ such that $E \subset \bigcup_{k \in K_j} U_{jk}$
and such that each restriction $s_j|_{U_{jk}}$
is a finite sum $\sum_{i \in I_{jk}} f_{jki}$
with $I_{jk} \subset I$, and $f_{jki}$ in the summand
$\mathcal{O}_X$ corresponding to $i \in I$. Set
$I_j = \bigcup_{k \in K_j} I_{jk}$. This is a finite set.
Since $U \subset E \subset \bigcup_{k \in K_j} U_{jk}$
the section $s_j|_U$ is a section of the finite direct sum
$\bigoplus_{i \in I_j} \mathcal{O}_X$.
By Lemma \ref{lemma-limits-colimits}
we see that actually $s_j|_U$ is a sum
$\sum_{i \in I_j} f_{ij}$ and
$f_{ij} \in \mathcal{O}_X(U) = \Gamma(U, \mathcal{O}_U)$.

\medskip\noindent
At this point we can define a module $M$ as the cokernel of the map
$$
\bigoplus\nolimits_{j \in J}
\Gamma(U, \mathcal{O}_U)
\longrightarrow
\bigoplus\nolimits_{i \in I}
\Gamma(U, \mathcal{O}_U)
$$
with matrix given by the $(f_{ij})$. By construction (2) of
Lemma \ref{lemma-construct-quasi-coherent-sheaves} we see that
$\mathcal{F}_M$ has the same presentation as $\mathcal{F}|_U$
and therefore $\mathcal{F}_M \cong \mathcal{F}|_U$.
\end{proof}

\begin{example}
\label{example-quasi-coherent}
Let $X$ be countably many copies $L_1, L_2, L_3, \ldots$
of the real line all glued together at $0$; a fundamental
system of neighbourhoods of $0$ being the collection
$\{U_n\}_{n \in \mathbf{N}}$, with $U_n \cap L_i = (-1/n, 1/n)$.
Let $\mathcal{O}_X$ be the sheaf of continuous real valued functions.
Let $f : \mathbf{R} \to \mathbf{R}$ be a continuous function
which is identically zero on $(-1, 1)$ and identically $1$
on $(-\infty, -2) \cup (2, \infty)$. Denote $f_n$ the continuous
function on $X$ which is equal to $x \mapsto f(nx)$ on each
$L_j = \mathbf{R}$. Let $1_{L_j}$ be the characteristic function
of $L_j$. We consider the map
$$
\bigoplus\nolimits_{j \in \mathbf{N}}
\mathcal{O}_X
\longrightarrow
\bigoplus\nolimits_{j, i \in \mathbf{N}}
\mathcal{O}_X, \quad
e_j \longmapsto \sum\nolimits_{i \in \mathbf{N}} f_j 1_{L_i} e_{ij}
$$
with obvious notation. This makes sense because this sum is locally
finite as $f_j$ is zero in a neighbourhood of $0$. Over $U_n$
the image of $e_j$, for $j > 2n$ is not a finite linear combination
$\sum g_{ij} e_{ij}$ with $g_{ij}$ continuous. Thus there is
no neighbourhood of $0 \in X$ such that the displayed map is given by
a ``matrix'' as in the proof of Lemma \ref{lemma-quasi-coherent-module} above.

\medskip\noindent
Note that $\bigoplus\nolimits_{j \in \mathbf{N}} \mathcal{O}_X$
is the sheaf associated to the free module with basis $e_j$
and similarly for the other direct sum. Thus we see that a
morphism of sheaves associated to modules in general even
locally on $X$ does not come from a morphism of modules.
Similarly there should be an example of a ringed space $X$
and a quasi-coherent $\mathcal{O}_X$-module $\mathcal{F}$
such that $\mathcal{F}$ is not locally of the form $\mathcal{F}_M$.
(Please email if you find one.)
Moreover, there should be examples of locally compact spaces
$X$ and maps $\mathcal{F}_M \to \mathcal{F}_N$ which also do
not locally come from maps of modules (the proof of Lemma
\ref{lemma-quasi-coherent-module} shows this cannot happen
if $N$ is free).
\end{example}












